\documentclass[12pt]{article}

\usepackage[czech]{babel}
\usepackage[utf8]{inputenc}
\usepackage[T1]{fontenc}
\usepackage{lingmacros}
\usepackage{tree-dvips}
\usepackage{listings}
\usepackage{xstring}
\usepackage{xparse}
\usepackage{booktabs}
\usepackage{longtable}
\usepackage[hidelinks]{hyperref}
\usepackage{parskip}

\author{\textbf{Tomáš Tobrman}}
\title{Semestrální práce \textbf{KIV/ZOS}}
\date{}

\renewcommand*\contentsname{Obsah}

\begin{document}

\maketitle
\thispagestyle{empty}
\newpage

\tableofcontents
\thispagestyle{empty}
\newpage

\setcounter{page}{1}

\section{Úvod}

\subsection{Zadání}

Vytvoření konzolového rozhraní s~příkazy pro práci s~pseudosystémem založeným na souborovém systému FAT.

\subsection{Implementace}

Projekt byl implementován v programovacím jazyce C pro operační systémy Linux/Unix.
Každému příkazu připadá vlastní zdrojový soubor, ve kterém se nachází jeho funkcionalita.
Pole příkazů je pak definováno ve zdrojovém souboru \texttt{utils.c}, proti kterému se porovnává uživatelský vstup.

\section{Použití}\label{sec:usage}

Před spuštěním programu je třeba zahájit jeho kompilaci.
Na systémech Linux/Unix je možné tohoto docílit pomocí příkazu \texttt{make}.
Po úspěšném sestavení program lze spustit z~příkazové řádky ve složce \texttt{bin}.

\subsection{Příkazy}\label{sec:commands}

Tato sekce obsahuje popis a~příklady využití všech příkazů implementovaných v~konzolovém rozhraní.

\subsubsection{\texttt{cat}}\label{sec:cat}

Vypíše obsah souboru.

\texttt{\$ cat file}

Příkaz očekává vždy jeden argument, který je cestou k~souboru, jehož obsah chce uživatel vypsat.

\subsubsection{\texttt{cd}}\label{sec:cd}

Změní cestu aktuálního pracovního adresáře.

\texttt{\$ cd path}

\texttt{\$ cd \/}

\texttt{\$ cd}

Příkaz očekává až jeden argument, který odpovídá cestě do adresáře.
Pokud není poskytnuta cesta, pracovní adresář je změněn na domácí složku uživatele (v tomto případě kořenový adresář \texttt{\/}).

\subsubsection{\texttt{clear}}\label{sec:clear}

\subsubsection{\texttt{cp}}\label{sec:cp}

\subsubsection{\texttt{defrag}}\label{sec:defrag}

\subsubsection{\texttt{exit}}\label{sec:exit}

\subsubsection{\texttt{format}}\label{sec:format}

\subsubsection{\texttt{help}}\label{sec:help}

\subsubsection{\texttt{incp}}\label{sec:incp}

\subsubsection{\texttt{info}}\label{sec:info}

\subsubsection{\texttt{load}}\label{sec:load}

\subsubsection{\texttt{ls}}\label{sec:ls}

\subsubsection{\texttt{mkdir}}\label{sec:mkdir}

\subsubsection{\texttt{mv}}\label{sec:mv}

\subsubsection{\texttt{outcp}}\label{sec:outcp}

\subsubsection{\texttt{pwd}}\label{sec:pwd}

\subsubsection{\texttt{rm}}\label{sec:rm}

\subsubsection{\texttt{rmdir}}\label{sec:rmdir}

\section{Závěr}

\end{document}
